% BANKS-SOURCE: pdf=174 print=164 kind=source-page
% BANKS-SOURCE: pdf=174 print=164 kind=prose
infrared-free because it is well approximated by a perturbation expansion in
$\lambda(t)$ at low momentum, even if $\lambda$ is not small.

But there is something confusing and sick here. Remember, we defined a QFT by
tuning relevant parameters to zero as we took the cut-off to infinity. The RG
trajectory approaches very close to the scale-invariant fixed-point theory
(which in our case is massless free-field theory) and then runs away from it at
the renormalization scale $\mu$, which is tuned to remain finite as the cut-off
goes to infinity. If we look back along the RG trajectory, from the scale $\mu$
to the UV, we can see only the fixed point. So a renormalized theory defined by
perturbing the massless free fixed point must behave like the free-field theory
at very high momenta. We have found the renormalized coupling getting strong in
this limit, which is a contradiction.

The clue to this behavior is infrared freedom. $\lambda(t)$ is getting weak in
the IR. In other words, it is behaving like an irrelevant (marginally
irrelevant) coupling. Its value in the renormalized theory should be fixed by the
other couplings. But the only other coupling is the mass. We can solve the
massive free-field theory exactly, and we find no connected four-point function.
Therefore the renormalized $\phi^4$ coupling is zero (this is often referred to
as triviality of $\phi^4$ theory).

Why then were we able to find a finite perturbation expansion to all orders in
the renormalized coupling $\lambda$? The clue lies in noting that the place
where the coupling gets strong in the UV (it blows up there but we can no longer
trust the perturbation expansion at the singularity) is at $t$ of order
$1/\lambda$. If $\lambda$ is small this is a momentum scale exponentially higher
than the mass of the particles in the theory $O(m^2)$. So we can put a cut-off
in the theory at an exponentially large scale, and get perfectly predictive
physics by ignoring the cut-off, up to terms of order $\ee^{-c/\lambda}$. In
other words, theories with marginally irrelevant parameters are almost as good
as mathematically well-defined QFTs, as long as the renormalized values of the
marginally irrelevant couplings are small enough to be in the perturbative
regime. The point where the perturbative coupling blows up is called the Landau
pole. As long as one uses a renormalized perturbation expansion at energies way
below the Landau pole, one is on safe ground.

% BANKS-SOURCE: pdf=174 print=164 kind=section id=9.8
\subsection{Renormalization of QED at one loop}

% BANKS-SOURCE: pdf=174 print=164 kind=prose
As a further example, we will treat the renormalization of quantum
electrodynamics at one loop. Power counting at the Gaussian fixed point for
electrodynamics in any covariant gauge suggests that the interaction $A_\mu J^\mu$
is marginal when $J^\mu$ is taken to be the current of either spin-zero or
spin-$\tfrac12$ particles. The additional $A_\mu^2\phi^*\phi$ interaction
necessary to the gauge-invariant Lagrangian for charged spin-zero particles is
also marginal. The scalar- and fermion-mass terms are relevant, and the only
other gauge-invariant marginal operator is a quartic coupling
$(\phi^*\phi)^2$. If we have multiple scalars and fermions the only additional
marginal couplings for generic values of the charges of different fields would be
of the form $\lambda_{ij}(\phi_i^*\phi_i)(\phi_j^*\phi_j)$. For special values
of the charges we could also have Yukawa couplings, cubic couplings of the
scalars, or both. However, even when the charges allow these couplings, we can
invoke a $\phi_i\to-\phi_i$ symmetry to forbid them.

% BANKS-SOURCE: pdf=175 print=165 kind=source-page
% BANKS-SOURCE: pdf=175 print=165 kind=prose
Wilson has emphasized that general renormalization theory does not require us to
use a cut-off that preserves symmetries or gauge invariances (redundancies) of
the classical Lagrangian, at least as long as the Gaussian fixed-point theory has
only a finite number of relevant and marginal perturbations.\footnote[11]{Which
include all the interactions in the classical Lagrangian. It is the failure of
this condition that makes spontaneously broken non-abelian gauge theories
non-renormalizable in the unitary gauge.} We simply introduce any cut-off we
like and then tune the relevant and marginal parameters to make the quantum
theory satisfy the quantum Ward identities. This could fail only if there were
violations of the Ward identities, which cannot be removed by adding a local
counterterm to the action. Schwinger--Adler--Bell--Jackiw anomalies, which we
studied in Chapter 8, are the unique example of the latter phenomenon.

Nonetheless, retaining symmetries at every step greatly simplifies the
renormalization process. For simplicity then, we will insist on using a
gauge-invariant regulator, so that the most general marginal and relevant local
Euclidean Lagrangian has the form
% BANKS-SOURCE: pdf=175 print=165 kind=equation
\[
  \mathcal{L}=\frac{Z_3}{4}F_{\mu\nu}^2
  +\frac{1}{2}Z_s|D_\mu\phi|^2
  +Z_f\bar\psi(\ii\gamma^\mu D_\mu-m_0)\psi
  +\mu_0^2\phi_i^*\phi_i
  +\frac{e_0^2v_0^{(4)}}{2}(\phi_i^*\phi_i)^2
  +\frac{1}{2\kappa_0}(\partial_\mu A_\mu)^2.
\]
The displayed Lagrangian and the loop formulae below are Euclidean. There
$\gamma^\mu$ denotes a continued matrix with
$\{\gamma^\mu,\gamma^\nu\}=2\delta^{\mu\nu}$, and $d_{\mathrm{reg}}$ is the
continued integration dimension near the physical $D=4$, with
$\epsilon\equiv4-d_{\mathrm{reg}}$ for pole notation. The Lorentzian
mostly-plus convention is $\{\gamma^\mu,\gamma^\nu\}=-2\eta^{\mu\nu}$,
$\slashed p=\gamma^\mu p_\mu$, with Dirac operator
$-\ii\slashed\partial-m$ and propagator
$S_F(p)=-\ii(\slashed p+m)/(p^2+m^2-\ii0)$. For any charged field with charge $q$, the covariant derivative is
$D_\mu\equiv\partial_\mu-\ii e_0qA_\mu$. We have chosen the loop-counting
parameter $e_0$ to be such that $e_0^2/(4\pi)$ is the bare fine-structure
constant, and are already working with fields that, at tree level, have
canonically normalized kinetic terms. In tree approximation, all the $Z$ factors
are equal to $1$, particle masses are given by $m=m_0$, $\mu=\mu_0$, and the
quartic coupling is $e_0^2v_0^{(4)}/2$.

General renormalization theory leads us to believe that we can eliminate all
cut-off dependence in physical amplitudes by tuning the $Z$ factors and the bare
parameters (those subscripted with a zero) as the cut-off is taken to infinity.
We will attempt to show below that the divergent part of the one-loop quantum
action, $\Gamma[A,\psi,\phi]$, has precisely the form given above, so that the
divergences can be eliminated as claimed. We will demonstrate this below, to
one-loop order, for a single spinor field. The generalization to the case of
multiple charged scalar and spinor fields should be a straightforward exercise
for the diligent reader.

\BanksImplicitHook{I-CH09-010}

For photons, there is a gauge-invariant generalization of the sort of
regularization procedure we introduced for scalars:
% BANKS-SOURCE: pdf=175 print=165 kind=equation
\[
  F_{\mu\nu}^2\longrightarrow
  F_{\mu\nu}K\!\left(\frac{\partial^2}{\Lambda^2}\right)F_{\mu\nu}.
\]
However, if we attempt to generalize this to charged fields, retaining gauge
invariance, we introduce an infinite number of apparently irrelevant
interactions. There are only three known classes of gauge-invariant
regularization prescription. The first is lattice gauge theory, but we will
insist on retaining Lorentz invariance. The second class uses

% BANKS-SOURCE: pdf=176 print=166 kind=source-page
the fact that charged fields appear only quadratically in the QED Lagrangian.\footnote[12]{We
can eliminate the quartic scalar interaction by introducing an auxiliary field
$g_0^2v_0^{(4)}(\phi^*\phi)^2/2\to\sigma^2/2+\ii\sigma
g_0\sqrt{v_0^{(4)}}\,\phi^*\phi$.} We can
formally integrate them out, obtaining a power of a functional determinant of a
gauge-covariant operator $D^2$ or $\gamma^\mu D_\mu$. There are many ways to
regulate the product over the gauge-invariant eigenvalues of this operator. For
example, Pauli and Villars pointed out that
% BANKS-SOURCE: pdf=176 print=166 kind=equation
\[
  \prod_i\det(D+\mu_i^{p_i})^{\epsilon_i}
\]
was finite for finite values of $\mu_i$ and appropriate choices of
$\epsilon_i=\pm1$. Here $D=-D^2$ for spin zero (where $p_i=2$) and
$\ii\gamma^\mu D_\mu$ (where $p_i=1$ and $\mu_i=\ii M_i$). We need only a
finite number of extra Pauli--Villars fields, but some of them must have wrong
statistics ($\epsilon_i=1$ for spin zero or $-1$ for spin $\tfrac12$). The
masses of the extra fields act as regulators, and are taken to infinity at the
end of the calculation.

Schwinger's gauge-invariant proper-time formula uses the observation that
% BANKS-SOURCE: pdf=176 print=166 kind=equation
\[
  \ln\det\left[\frac{-D^2+m^2}{-\partial^2+m^2}\right]
  =\int_0^\infty\frac{\dd t}{t}
  \left[\ee^{-t(-D^2+m^2)}-\ee^{-t(-\partial^2+m^2)}\right].
\]
If we cut the lower limit of the $t$ integral off at $t_0$ we suppress the
contribution of large eigenvalues of $-D^2$ and get a finite, gauge-invariant
functional of $A_\mu$. $t_0$ has dimensions of inverse mass squared and defines a
cut-off. In order to use this formula in spinor electrodynamics, we must
recognize that the determinant of the Euclidean Dirac operator is formally
positive and independent of the sign of the fermion mass, and so is equal to the
positive square root of the determinant of
$(-D^2+m^2+\ii\sigma^{\mu\nu}F_{\mu\nu}/2)$. We can then apply proper-time
regularization to the latter quantity.

The combination of either of these two methods with the higher-derivative
cut-off of the photon propagator gives us a finite, gauge-invariant,
Lorentz-invariant answer for all Green functions in spinor and scalar QED. The
answer does not obey the quantum requirement of unitarity until we take the
cut-off to infinity. Unfortunately, neither of these methods generalizes to
non-abelian gauge theory. The only known Lorentz-invariant, gauge-invariant
regulator for the non-abelian case is dimensional regularization (DR), and this
is the method we will use. We begin by studying the wave-function
renormalization of the photon.

Figure~\ref{fig:9.6} shows the 1PI photon two-point-function graph at one loop in spinor
QED.
% BANKS-SOURCE: pdf=176 print=166 kind=figure id=9.6
\begin{figure}[htbp]
  \centering
  % BANKS-SOURCE: pdf=176 print=166 kind=figure id=9.6
\begin{tikzpicture}[x=1cm,y=1cm,baseline=-.4ex]
  % fermion loop, counter-clockwise (arrow left on top, right on bottom)
  \draw[line width=.55pt,
        decoration={markings,
          mark=at position 0.25 with {\arrow{Latex[length=2.4mm,width=2.4mm]}},
          mark=at position 0.75 with {\arrow{Latex[length=2.4mm,width=2.4mm]}}},
        postaction={decorate}]
    (0,0) circle (1.0);
  % incoming and outgoing photons
  \draw[decorate,decoration={snake,amplitude=2.2pt,segment length=6.5pt},line width=.55pt]
    (-2.55,0)--(-1.0,0);
  \draw[decorate,decoration={snake,amplitude=2.2pt,segment length=6.5pt},line width=.55pt]
    (1.0,0)--(2.55,0);
  % vertices
  \fill (-1.0,0) circle (2.1pt);
  \fill (1.0,0) circle (2.1pt);
\end{tikzpicture}

  \caption{Vacuum polarization in QED.}
  \label{fig:9.6}
\end{figure}

% BANKS-SOURCE: pdf=177 print=167 kind=source-page
The value of the diagram is
% BANKS-SOURCE: pdf=177 print=167 kind=equation
\[
  \Pi_{\alpha\beta}(p)=e_0^2\int\frac{\dd^{d_{\mathrm{reg}}}q}{(2\pi)^{d_{\mathrm{reg}}}}
  \operatorname{tr}\!\left[
  \frac{\gamma_\alpha(\gamma^\mu q_\mu-\ii m_0)\gamma_\beta
  (\gamma^\mu(p+q)_\mu-\ii m_0)}
  {(q^2+m_0^2)[(p+q)^2+m_0^2]}\right].
\]
Our first task is to show that this formula is transverse:
% BANKS-SOURCE: pdf=177 print=167 kind=equation
\[
  \Pi_{\alpha\beta}(p)=\Pi(p^2)
  (\delta_{\alpha\beta}p^2-p_\alpha p_\beta).
\]
The less-than-alert reader will have missed the fact that, in our discussion of
the gauge-invariant quantum action above, we did not allow for a one-loop
correction to the gauge-fixing parameter $\kappa_0$. The fact that the 1PI photon
propagator is transverse (to all orders) guarantees this. So we compute
% BANKS-SOURCE: pdf=177 print=167 kind=equation
\[
  p^\alpha\Pi_{\alpha\beta}\propto e_0^2\int\frac{\dd^{d_{\mathrm{reg}}}q}{(2\pi)^{d_{\mathrm{reg}}}}
  \operatorname{tr}\!\left[
  \frac{p^\alpha\gamma_\alpha(\gamma^\mu q_\mu-\ii m_0)\gamma_\beta
  (\gamma^\mu(p+q)_\mu-\ii m_0)}
  {(q^2+m_0^2)[(p+q)^2+m_0^2]}\right].
\]
Write
% BANKS-SOURCE: pdf=177 print=167 kind=equation
\[
  p^\alpha\gamma_\alpha
  =[(p+q)^\alpha\gamma_\alpha-\ii m_0]
  -(q^\alpha\gamma_\alpha-\ii m_0),
\]
in order to obtain
% BANKS-SOURCE: pdf=177 print=167 kind=equation
\[
  p^\alpha\Pi_{\alpha\beta}\propto e_0^2\int\frac{\dd^{d_{\mathrm{reg}}}q}{(2\pi)^{d_{\mathrm{reg}}}}
  \operatorname{tr}\!\left[
  \frac{q^\alpha\gamma_\alpha-\ii m_0}{q^2+m_0^2}
  -\frac{(p+q)^\alpha\gamma_\alpha-\ii m_0}{(p+q)^2+m_0^2}
  \right].
\]
DR is defined so that we can shift variables of integration, so
$p^\alpha\Pi_{\alpha\beta}=0$. Now take the trace $\Pi^\mu{}_{\mu}$, and use
transversality to obtain
\BanksImplicitHook{I-CH09-011}

% BANKS-SOURCE: pdf=177 print=167 kind=equation
\[
  (d_{\mathrm{reg}}-1)p^2\Pi(p^2)=e_0^2\int\frac{\dd^{d_{\mathrm{reg}}}q}{(2\pi)^{d_{\mathrm{reg}}}}
  \operatorname{tr}\!\left[
  \frac{\gamma_\alpha(\gamma^\mu q_\mu-\ii m_0)\gamma^\alpha
  (\gamma^\mu(p+q)_\mu-\ii m_0)}
  {(q^2+m_0^2)[(p+q)^2+m_0^2]}\right].
\]
We use the identities
% BANKS-SOURCE: pdf=177 print=167 kind=equation
\[
  \gamma^\alpha\gamma^\mu q_\mu\gamma_\alpha=(2-d_{\mathrm{reg}})\gamma^\mu q_\mu,
  \qquad
  \gamma^\alpha m_0\gamma_\alpha=d_{\mathrm{reg}}m_0,
\]
to rewrite this as
% BANKS-SOURCE: pdf=177 print=167 kind=equation
\[
  (d_{\mathrm{reg}}-1)p^2\Pi(p^2)=\frac{e_0^2d_{\mathrm{spin}}}{(2\pi)^{d_{\mathrm{reg}}}}
  \int\dd^{d_{\mathrm{reg}}}q\,
  \frac{q\mathbin{\cdot}(p+q)(2-d_{\mathrm{reg}})-d_{\mathrm{reg}}m_0^2}
  {(q^2+m_0^2)[(p+q)^2+m_0^2]}.
\]
Our next piece of trickery is
% BANKS-SOURCE: pdf=177 print=167 kind=equation
\[
  q\mathbin{\cdot}(p+q)=\frac{1}{2}(p+q)^2-\frac{1}{2}q^2-\frac{1}{2}p^2+q^2.
\]
The first two terms cancel out after integration and shifting of variables. So
we obtain
% BANKS-SOURCE: pdf=177 print=167 kind=equation
\[
  (d_{\mathrm{reg}}-1)p^2\Pi(p^2)=\frac{e_0^2d_{\mathrm{spin}}}{(2\pi)^{d_{\mathrm{reg}}}}
  \int\dd^{d_{\mathrm{reg}}}q\,
  \frac{[(d_{\mathrm{reg}}-2)/2]p^2-d_{\mathrm{reg}}m_0^2+(2-d_{\mathrm{reg}})q^2}
  {(q^2+m_0^2)[(p+q)^2+m_0^2]}.
\]

% BANKS-SOURCE: pdf=178 print=168 kind=source-page
% BANKS-SOURCE: pdf=178 print=168 kind=prose
To evaluate the momentum integrals, we use Schwinger's parametric formula
$1/A=\int_0^\infty\dd\alpha\,\ee^{-\alpha A}$ to get
% BANKS-SOURCE: pdf=178 print=168 kind=equation
\[
\begin{aligned}
 (d_{\mathrm{reg}}-1)p^2\Pi(p^2)={}&\frac{e_0^2d_{\mathrm{spin}}}{(2\pi)^{d_{\mathrm{reg}}}}
 \int_0^\infty\dd\alpha\,\dd\beta\,
 \ee^{-(\alpha+\beta)m_0^2}\ee^{-\beta p^2}
 \int\dd^{d_{\mathrm{reg}}}q\,\ee^{-(\alpha+\beta)q^2}\\[-2pt]
 &\quad\times\left[\frac{d_{\mathrm{reg}}-2}{2}p^2-d_{\mathrm{reg}}m_0^2
 +(2-d_{\mathrm{reg}})\frac{\nabla_p^2}{4\beta^2}\right]
 \ee^{2\beta p\mathbin{\cdot}q}\\[2pt]
={}&\frac{e_0^2d_{\mathrm{spin}}}{(2\sqrt\pi)^{d_{\mathrm{reg}}}}
 \int_0^\infty\dd\alpha\,\dd\beta\,
 \frac{\ee^{-(\alpha+\beta)m_0^2}\ee^{-\beta p^2}}
 {(\alpha+\beta)^{d_{\mathrm{reg}}/2}}\\[-2pt]
 &\quad\times\left[\frac{d_{\mathrm{reg}}-2}{2}p^2-d_{\mathrm{reg}}m_0^2
 +(2-d_{\mathrm{reg}})\frac{\nabla_p^2}{4\beta^2}\right]
 \ee^{\beta^2p^2/(\alpha+\beta)}.
\end{aligned}
\]
In order to understand the last equality, the knowledgeable reader will have
used the formula for Euler's gamma function,
% BANKS-SOURCE: pdf=178 print=168 kind=equation
\[
  \Gamma(t)=\int\dd u\,u^{t-1}\ee^{-u},
\]
as well as the identities
% BANKS-SOURCE: pdf=178 print=168 kind=equation
\[
  \Gamma(1+t)=t\Gamma(t),\qquad
  \Gamma(t)=(t-1)!,\quad\text{if $t$ is an integer $\geq1$},
  \qquad
  \Gamma(t)\Gamma(1-t)=\frac{\pi}{\sin(\pi t)},
\]
from which it devolves that $\Gamma(1/2)=\sqrt\pi$.
\BanksImplicitHook{I-CH09-012}

It is now obvious that we should introduce $s\equiv\alpha+\beta$ and
$x\equiv\beta/(\alpha+\beta)$, and note that
$\dd\alpha\,\dd\beta=s\,\dd s\,\dd x$. We obtain the parametric formula
% BANKS-SOURCE: pdf=178 print=168 kind=equation
\[
\begin{aligned}
 (d_{\mathrm{reg}}-1)p^2\Pi(p^2)={}&\frac{e_0^2d_{\mathrm{spin}}}{(2\sqrt\pi)^{d_{\mathrm{reg}}}}
 \int_0^\infty\dd s\int_0^1\dd x\,
 \ee^{-s[m_0^2+x(1-x)p^2]}\\[-2pt]
 &\quad\times\left[\frac{d_{\mathrm{reg}}-2}{2}p^2-d_{\mathrm{reg}}m_0^2
 +(2-d_{\mathrm{reg}})\left(\frac{d_{\mathrm{reg}}}{2s}+x^2p^2\right)\right].
\end{aligned}
\]
Simple power counting gives a quadratic divergence in $\Pi_{\alpha\beta}$,
which would be momentum-independent and correspond to a photon mass
$\Lambda^2A_\mu^2$. The transverse formula for $\Pi_{\alpha\beta}$ shows that
such a term cannot be there. The only way we could get something that looks
like a photon mass would be from a pole in $\Pi(p^2)$ at $p^2=0$. This would
not be UV-divergent. In fact, in more than two space-time dimensions such a
pole does not appear (and even there it appears only if $m_0=0$). The only way
to get a massive photon is through the Higgs mechanism. In perturbation theory
this can happen only if there is a scalar field with a negative mass squared in
the tree-level Lagrangian. In that case we have to diagonalize the free
Lagrangian properly and

% BANKS-SOURCE: pdf=179 print=169 kind=source-page
we find only massive tree-level excitations. It must be, then, that the formula
we have written above vanishes when $p^2=0$. The diligent reader will verify
that fact. We thus obtain
% BANKS-SOURCE: pdf=179 print=169 kind=equation
\[
\begin{aligned}
 (d_{\mathrm{reg}}-1)\Pi(p^2)={}&\frac{e_0^2}{(2\sqrt\pi)^{d_{\mathrm{reg}}}}d_{\mathrm{spin}}
 \int_0^\infty\dd s\,s^{1-d_{\mathrm{reg}}/2}\int_0^1\dd x\\[-2pt]
 &\quad\times\left[d_{\mathrm{reg}}x(1-x)+(2-d_{\mathrm{reg}})\left(x^2-\frac12\right)
 \right]\ee^{-s[m_0^2+x(1-x)p^2]}.
\end{aligned}
\]
\BanksImplicitHook{I-CH09-013}

Poles when $d_{\mathrm{reg}}\to4$ will come from UV-divergent integrals at
$d_{\mathrm{reg}}=4$.
Ultraviolet divergences correspond to the small-$s$ region of integration. It
is now obvious that terms in our formula of order higher than $p^2$ in an
expansion about $0$ will not have divergences in the small-$s$ region. Thus,
the divergence resides in $\Pi(0)$, which means that it has the form of a
correction to the tree-level relation $Z_3=1$. In fact
\BanksImplicitHook{I-CH09-014}

% BANKS-SOURCE: pdf=179 print=169 kind=equation
\[
  (d_{\mathrm{reg}}-1)\Pi(0)=\frac{e_0^2}{(2\sqrt\pi)^{d_{\mathrm{reg}}}}
  d_{\mathrm{spin}}m_0^{d_{\mathrm{reg}}-4}\Gamma\!\left(2-\frac{d_{\mathrm{reg}}}{2}\right).
\]
The peculiar factor of $m_0^{d_{\mathrm{reg}}-4}$ in this formula points up something
important. Away from four dimensions, the electromagnetic coupling is not
dimensionless. DR introduces no explicit cut-off with the dimensions of mass,
so the dimensions in formulae must be made up by powers of $m_0$ or some other
mass parameter.

The pole in the 1PI photon Green function can all be attributed to the $Z_3$
factor, but there is an ambiguity regarding how much finite part to keep when
we define $Z_3$. DR allows us to keep ``just the pole,'' but when we express
the answer in terms of $e_0$ this appears to introduce a dependence on $m_0$
because of dimensional analysis. However, we can introduce an arbitrary mass
scale $\mu$ in place of $m_0$ to define the split between $Z_3$ and the
renormalized field in the formula
% BANKS-SOURCE: pdf=179 print=169 kind=equation
\[
  A_\mu=Z_3^{1/2}A_\mu^R.
\]
We thus write
% BANKS-SOURCE: pdf=179 print=169 kind=equation
\[
  (Z_3-1)=\frac{2}{3}\frac{\alpha_0}{\pi\mu^{4-d_{\mathrm{reg}}}}
  \frac{1}{d_{\mathrm{reg}}-4},
\]
where we have evaluated the residue of the pole at $d_{\mathrm{reg}}=4$. This prescription
for parametrizing the renormalized theory is called minimal subtraction (MS).
Note that
% BANKS-SOURCE: pdf=179 print=169 kind=equation
\[
  \Gamma\!\left(2-\frac{d_{\mathrm{reg}}}{2}\right)
  \longrightarrow\frac{1}{2-d_{\mathrm{reg}}/2}-\gamma,
\]
where $\gamma$ is the Euler--Mascheroni constant. There is a modified
prescription, called $\overline{\mathrm{MS}}$, which keeps certain finite parts
in addition to the pole, and removes many of the $\gamma$ factors which would
otherwise appear in renormalized formulae (see Peskin and Schroeder
\cite{banks-ref-33} for details).

% BANKS-SOURCE: pdf=180 print=170 kind=section id=9.8.1
\subsection{Renormalization of the fermion propagator and the vertex}

% BANKS-SOURCE: pdf=180 print=170 kind=prose
We now turn to the other two divergent one-loop diagrams in spinor QED, the
fermion self-energy and the photon fermion vertex. The gauge-invariant form of
the Lagrangian suggests that these are renormalized by the same multiplicative
factor. The 1PI vertex function $\Gamma^\mu(p,p+q)$ is related to the Fourier
transform of the Green function of the electromagnetic current,
% BANKS-SOURCE: pdf=180 print=170 kind=equation
\[
  \langle J^\mu(x)\psi(y)\bar\psi(z)\rangle,
\]
by multiplying it by inverse fermion propagators on the external legs and by
$e_0$, the bare charge. The current Green function satisfies the Ward identity
% BANKS-SOURCE: pdf=180 print=170 kind=equation
\[
  \partial_\mu^x\langle J^\mu(x)\psi(y)\bar\psi(z)\rangle
  =\ii\delta^{d_{\mathrm{reg}}}(x-y)\langle\psi(y)\bar\psi(z)\rangle
  -\ii\delta^{d_{\mathrm{reg}}}(x-z)\langle\psi(y)\bar\psi(z)\rangle.
\]
Written in terms of the vertex function, this identity is
% BANKS-SOURCE: pdf=180 print=170 kind=equation
\[
  q_\mu\Gamma^\mu(p,p+q)=e_0\left[\gamma^\mu q_\mu
  +\Sigma(p+q)-\Sigma(p)\right],
\]
where $\Sigma(p)$ is the sum of all 1PI self-energy graphs (the notation here
differs a bit from standard texts because I have included the tree-level term in
the definition of $\Gamma_\mu$).

The vertex function itself is dimensionless, so the one-loop integral defining it
is at most logarithmically divergent. Consequently, all derivatives of
$\Gamma_\mu$ with respect to $q$ are given by finite expressions in four
dimensions and will not have poles at $d_{\mathrm{reg}}=4$. The divergent part of
$\Gamma_\mu$
is thus independent of $q$ at one loop. A similar argument shows that it is
independent of $p$ as well. Reflection invariance of QED shows that it cannot
involve $\gamma_5$, so we must have
% BANKS-SOURCE: pdf=180 print=170 kind=equation
\[
  \Gamma^\infty_\mu=e_0(Z_1-1)\gamma_\mu.
\]
The Ward identity shows that this must be related to a divergent term in
$\Sigma^\infty(p)=(Z_2-1)\gamma_\mu p^\mu$, with $Z_1=Z_2$.

We therefore turn to the self-energy diagram of Figure~\ref{fig:9.7} to evaluate $Z_2$.
The result will depend on the gauge parameter $\kappa_0$ in the free-photon
propagator
% BANKS-SOURCE: pdf=180 print=170 kind=equation
\[
  \frac{\delta_{\mu\nu}-q_\mu q_\nu/q^2}{q^2}
  +\kappa_0\frac{q_\mu q_\nu}{q^4}.
\]
We will notice that Landau (also called Lorentz) gauge, $\kappa_0=0$, has
certain nice features (it avoids some infrared divergences in low orders). One
way of understanding why the Landau-gauge fermion propagator is so nice is to
note that it is actually the value of a

% BANKS-SOURCE: pdf=180 print=170 kind=figure id=9.7
\begin{figure}[htbp]
  \centering
  % BANKS-SOURCE: pdf=180 print=170 kind=figure id=9.7
\begin{tikzpicture}[x=1cm,y=1cm,baseline=-.4ex]
  % fermion line with a mid-line arrow
  \draw[line width=.55pt,
        decoration={markings,
          mark=at position 0.52 with {\arrow{Latex[length=2.4mm,width=2.4mm]}}},
        postaction={decorate}]
    (-2.2,0)--(2.2,0);
  % photon arc from the left vertex over to the right vertex
  \draw[decorate,decoration={snake,amplitude=3.2pt,segment length=13pt},line width=.55pt]
    (-1.05,0) arc (180:0:1.05);
  \node[anchor=east] at (-2.3,.05) {$p$};
  \node[anchor=north] at (0,-.1) {$p+q$};
  \node[anchor=south] at (0,1.16) {$q$};
\end{tikzpicture}

  \caption{One-loop self-energy in QED.}
  \label{fig:9.7}
\end{figure}

% BANKS-SOURCE: pdf=181 print=171 kind=source-page
fairly simple non-local gauge-invariant operator, evaluated in Landau gauge.
Indeed, if we multiply the bilinear $\psi(x)\bar\psi(y)$ by
% BANKS-SOURCE: pdf=181 print=171 kind=equation
\[
  \ee^{\ii\int A_\mu(x)j^\mu(x)},
\]
where $\partial_\mu j^\mu(z)=\delta^{d_{\mathrm{reg}}}(z-x)-\delta^{d_{\mathrm{reg}}}(z-y)$, then we get a
gauge-invariant operator. On choosing $j^\mu=\partial^\mu A$, and noting that
the resulting $A$ vanishes at infinity (it is the four-dimensional analog of
the field of a dipole), we see, upon integrating by parts, that the non-local
exponential vanishes in Landau gauge.

We will do the calculation in Euclidean space. We write the one-loop
contribution to $\Sigma(p)$ as
% BANKS-SOURCE: pdf=181 print=171 kind=equation id=9.1
\begin{equation}
  \Sigma(p)=-\frac{e_0^2}{(2\pi)^{d_{\mathrm{reg}}}}\int\dd^{d_{\mathrm{reg}}}q\,
  \frac{\gamma_\alpha(\slashed p-\slashed q-\ii m_0)\gamma_\beta}
  {q^2[(p-q)^2+m_0^2]}
  \left(\delta^{\alpha\beta}-(1-\kappa_0)\frac{q^\alpha q^\beta}{q^2}\right).
  \label{eq:9.1}
  \tag{9.1}
\end{equation}
We introduce Schwinger parameters to write the denominator as
% BANKS-SOURCE: pdf=181 print=171 kind=equation
\[
  \int_0^\infty\dd\alpha\,\dd\beta\,
  \ee^{-(\alpha+\beta)q^2}\ee^{-\beta(p^2+m_0^2)}
  \ee^{2\beta p\mathbin{\cdot}q}.
\]
Note that the extra inverse power of $q^2$ in the last term of the numerator
turns into an extra power of $\alpha$ in the numerator of the parametrized
integral. To do the Gaussian integral over loop momentum, we shift the
integration variable to $r=q-xp$. Here we have introduced the standard one-loop
passage from Schwinger to Feynman parameters: $s=\alpha+\beta$,
$\beta=sx$, $\dd\alpha\,\dd\beta=s\,\dd s\,\dd x$. When we do the integrals,
terms linear in $r$ will integrate to zero, and we need to use
% BANKS-SOURCE: pdf=181 print=171 kind=equation
\[
  \int\dd^{d_{\mathrm{reg}}}r\,\ee^{-sr^2}=\left(\frac{\pi}{s}\right)^{d_{\mathrm{reg}}/2}
\]
and
% BANKS-SOURCE: pdf=181 print=171 kind=equation
\[
  \int\dd^{d_{\mathrm{reg}}}r\,\ee^{-sr^2}r_\mu r_\nu
  =\frac{\delta_{\mu\nu}}{2s}\left(\frac{\pi}{s}\right)^{d_{\mathrm{reg}}/2}.
\]
Keeping only terms even in $r$, the numerator is
% BANKS-SOURCE: pdf=181 print=171 kind=equation
\[
\begin{aligned}
&\gamma^\alpha(\slashed p(1-x)-\ii m_0)\gamma^\beta
 \left[\delta_{\alpha\beta}-(1-\kappa)(1-x)s
 (r_\alpha r_\beta+x^2p_\alpha p_\beta)\right]\\
&\qquad+\gamma^\alpha\slashed r\gamma^\beta x(1-\kappa)s
 (r_\alpha p_\beta+r_\beta p_\alpha).
\end{aligned}
\]
Note that we have inserted the renormalized gauge parameter rather than the
bare one, because the correction is of higher order in $e^2$. After doing the
$r$ integration, terms quadratic in $r$ acquire an extra power of $s$. Terms
proportional to $s^{1-d_{\mathrm{reg}}/2}$ give
$\Gamma(2-d_{\mathrm{reg}}/2)\sim2/(4-d_{\mathrm{reg}})$, but terms with
higher powers of $s$ give convergent integrals. Thus, the divergent part of the
self-energy is given by
% BANKS-SOURCE: pdf=181 print=171 kind=equation
\[
\begin{aligned}
  \Sigma^\infty={}&\frac{2}{4-d_{\mathrm{reg}}}\int_0^1\dd x\,
  \gamma^\mu(\slashed p(1-x)-\ii m_0)\gamma_\mu
  \left(1-\frac{(1-\kappa)(1-x)}{2}\right)\\
  &\qquad+d_{\mathrm{reg}}\slashed p\,x(1-x)(1-\kappa).
\end{aligned}
\]

% BANKS-SOURCE: pdf=182 print=172 kind=source-page
Finally, we use contraction identities for Dirac matrices and do the $x$
integral to obtain
% BANKS-SOURCE: pdf=182 print=172 kind=equation
\[
  \Sigma^\infty=-\frac{e^2}{8\pi^2(4-d_{\mathrm{reg}})}
  \left[\kappa(\slashed p-\ii m)+3\ii m\right].
\]
This can obviously be removed by rescaling the fields and renormalizing the
mass according to
% BANKS-SOURCE: pdf=182 print=172 kind=equation
\[
  Z_2=1+\kappa\frac{\alpha}{2\pi\epsilon}
\]
and
% BANKS-SOURCE: pdf=182 print=172 kind=equation
\[
  \delta m=-3\frac{\alpha}{2\pi\epsilon}m_0.
\]
\BanksImplicitHook{I-CH09-015}

There are several interesting points about these formulae. First, we find that
$\delta m\propto m_0$, so the mass renormalization is multiplicative and
vanishes if $m_0$ vanishes. This is a consequence of the extra chiral symmetry
of the massless theory. Although DR doesn't really preserve conservation of the
associated Noether current (this is the chiral anomaly we have discussed above),
it does conserve the charge in perturbation theory. As a consequence, although
$m_0$ has dimensions of mass, it does not have the sensitivity to the cut-off
scale we expect for a relevant parameter. Note also that $\delta m$ is
independent of $\kappa$.

DR has the peculiar property that
% BANKS-SOURCE: pdf=182 print=172 kind=equation
\[
  \int\dd^{d_{\mathrm{reg}}}p\,\frac{1}{p^2+M^2}\propto M^{d_{\mathrm{reg}}-2}
\]
and thus vanishes for $M=0$. That is, quadratically divergent integrals with
only massless propagators vanish in DR, but they do depend quadratically on the
masses of very heavy particles. For scalar fields, this leads to renormalizations
of any mass term proportional to $\Lambda^2$, coming from integrating out
particles of mass near the cut-off. On dimensional grounds we might have
expected the same for fermion masses. However, the chiral symmetry of massless
fermion systems guarantees that the corresponding massive systems have only
logarithmically divergent mass renormalization. We will return briefly to this
point when we discuss the concept of \emph{technical naturalness} below.

The second interesting point is that the fermion wave-function renormalization
is $\kappa$-dependent and vanishes at $\kappa=0$. In Chapter 6 we noted that
QED has IR divergences associated with massless-photon emission in the
scattering of charged particles. One can show using the RG equations of the
next section and the IR freedom of QED that, in leading-order RG-improved
perturbation theory, the fermion propagator has a cut rather than a pole at the
physical mass, with a gauge-dependent power law. This is problematic for the LSZ
formula for the S-matrix. In Landau gauge, the cut becomes a pole in this
leading-order approximation. The gauge-invariant operator which is equal to the
fermion field in Landau gauge includes the Coulomb field of charged particles,
the IR effect that is of leading order in perturbation theory. In higher orders,
we must include appropriate coherent states of transverse photons in the
definition of charged-particle states, in order to account for bremsstrahlung
radiation and cancel out the IR divergences.
\BanksImplicitHook{I-CH09-016}
